\documentclass[12pt,a4paper]{article}
\usepackage{ibnp-base}
\usepackage{booktabs}
\usepackage{tabularx}

\ibnptitle{Modelo Padrão de Projeto de Impacto Social}
\ibnpsubject{Plano de Trabalho e Proposta Oficial}
\ibnplocation{Guapó - GO}
\ibnpdate{\today}

\begin{document}

% ---------------------------------------------------------
% Cabeçalho / Título de Abertura
% ---------------------------------------------------------
\begin{center}
    {\color{ibnpNavy}\sffamily\LARGE\bfseries PROJETO DE IMPACTO SOCIAL}\\[0.3cm]
    {\color{ibnpBlue}\sffamily\large Plano de Trabalho e Proposta Oficial de Execução}\\[0.2cm]
    {\small\textbf{Igreja Batista Nacional da Paz de Guapó (IBNP)} \textbullet{} Guapó - GO}
\end{center}

\vspace{0.5cm}

% ---------------------------------------------------------
% 1. Identificação do Proponente
% ---------------------------------------------------------
\section{Identificação do Proponente}

A proponente é uma organização religiosa e beneficente sem fins lucrativos, com sede própria e atuação consolidada no município de Guapó, Estado de Goiás.

\begin{table}[h!]
\centering
\small
\renewcommand{\arraystretch}{1.3}
\begin{tabularx}{\textwidth}{@{}>{\bfseries\color{ibnpNavy}}p{4.2cm}X@{}}
\toprule
Razão Social Oficial: & Igreja Batista Nacional da Paz de Guapó \\
Nome Fantasia / Sigla: & IBNP Guapó \\
CNPJ: & 02.930.019/0001-62 \\
Data de Fundação: & 14 de janeiro de 1999 \\
Natureza Jurídica: & 322-0 - Organização Religiosa Sem Fins Lucrativos \\
Sede e Endereço Físico: & Rua Presidente Kennedy, Qd. 21, Lt. 13, Centro, Guapó – GO, CEP 75350-000 \\
Telefone / WhatsApp: & (62) 9870-0089 \\
E-mail Institucional: & \href{mailto:contato@ibnpguapo.org.br}{contato@ibnpguapo.org.br} \\
Portal Oficial: & \href{https://ibnpguapo.org.br}{https://ibnpguapo.org.br} \\
Registro Notarial: & 2º Serviço Notarial da Comarca de Guapó - GO \\
Filiação Denominacional: & Convenção Batista Nacional (CBN) e ORMIBAN Goiás \\
\bottomrule
\end{tabularx}
\end{table}

\vspace{0.3cm}

% ---------------------------------------------------------
% 2. Apresentação e Diagnóstico Sociocultural do Município
% ---------------------------------------------------------
\section{Apresentação e Diagnóstico Sociocultural}

O município de Guapó, localizado na Região Metropolitana de Goiânia, apresenta um contingente expressivo de famílias que demandam opções descentralizadas de desenvolvimento sociocultural, musical e esportivo no contraturno escolar. O diagnóstico local evidencia que a escassez de atividades formativas acessíveis eleva a vulnerabilidade de crianças e adolescentes.

A Igreja Batista Nacional da Paz de Guapó mantém histórico de dedicação e prestação de serviços à comunidade através da \textbf{Escola Social de Guapó}, iniciativa que congrega acolhimento comunitário, ações socioeducativas, arrecadação de donativos e programas culturais contínuos.

\ibnpheaderbox{Síntese do Diagnóstico Municipal}{
O presente plano de trabalho destina-se a atender carências socioculturais em Guapó-GO, oferecendo estrutura formativa, acompanhamento integral e metodologia participativa totalmente gratuita para a comunidade local.
}

\vspace{0.3cm}

% ---------------------------------------------------------
% 3. Justificativa de Impacto Social
% ---------------------------------------------------------
\section{Justificativa de Impacto Social}

A realização de projetos socioculturais constitui ferramenta comprovada de promoção da cidadania, integração comunitária, prevenção à ociosidade e fortalecimento dos vínculos familiares. Ao proporcionar acesso orientado a práticas pedagógicas e artísticas, o projeto estimula competências socioemocionais fundamentais, tais como cooperação, disciplina, comunicação e autoestima.

Ademais, a iniciativa dialoga diretamente com as diretrizes do Estatuto da Criança e do Adolescente (ECA) e com as políticas públicas de fomento à cultura e proteção social básica, garantindo atendimento humanizado e sem qualquer discriminação.

\vspace{0.3cm}

% ---------------------------------------------------------
% 4. Objetivos Geral e Específicos
% ---------------------------------------------------------
\section{Objetivos Geral e Específicos}

\subsection{Objetivo Geral}
Promover o desenvolvimento sociocultural de crianças, adolescentes e famílias residentes no município de Guapó-GO, por meio de oficinas formativas continuadas, proporcionando vivência coletiva e integração comunitária.

\subsection{Objetivos Específicos}
\begin{enumerate}
    \item Ofertar vagas inteiramente gratuitas com atendimento presencial e suporte metodológico adequado.
    \item Desenvolver habilidades cognitivas, artísticas e expressivas através de módulos teóricos e práticos.
    \item Estimular o fortalecimento da convivência familiar e comunitária através de apresentações públicas de encerramento.
    \item Monitorar e certificar 100\% dos participantes com frequência regular ao término das atividades.
\end{enumerate}

\vspace{0.3cm}

% ---------------------------------------------------------
% 5. Metodologia de Execução e Público Atendido
% ---------------------------------------------------------
\section{Metodologia de Execução e Público Atendido}

A metodologia fundamenta-se na aprendizagem ativa, inclusiva e lúdica, distribuída em módulos formativos com dinâmicas de grupo, exercícios práticos e ensaios coletivos.

\begin{itemize}
    \item \textbf{Público-Alvo:} Crianças e adolescentes da rede pública de ensino e comunidade em situação de vulnerabilidade de Guapó-GO.
    \item \textbf{Capacidade de Atendimento:} [Número de Beneficiários Diretos] participantes por turma.
    \item \textbf{Carga Horária Total:} [Carga Horária] horas distribuídas em encontros regulares.
    \item \textbf{Critérios de Seleção:} Inscrição voluntária aberta, ordem de procura, prioridade para famílias em situação de vulnerabilidade cadastrada e termo de autorização dos responsáveis legais.
    \item \textbf{Gratuidade:} Garantia de 100\% de gratuidade no fornecimento de material didático, alimentação e certificação.
\end{itemize}

\newpage

% ---------------------------------------------------------
% 6. Cronograma Físico de Atividades
% ---------------------------------------------------------
\section{Cronograma Físico de Atividades}

O cronograma físico divide-se em 5 etapas estratégicas estruturadas para assegurar a excelência na execução e o acompanhamento transparente:

\begin{table}[h!]
\centering
\small
\renewcommand{\arraystretch}{1.2}
\begin{tabularx}{\textwidth}{@{}>{\bfseries}l X c c@{}}
\toprule
Etapa & Atividade / Meta Física & Período & Responsável \\
\midrule
1. Planejamento & Alinhamento pedagógico, materiais e logística & Mês 1 & Coordenação \\
2. Mobilização & Divulgação pública, abertura e validação de inscrições & Mês 1 & Equipe IBNP \\
3. Execução & Realização dos encontros presenciais e oficinas formativas & Mês 2 & Instrutores \\
4. Culminância & Apresentação pública comunitária e entrega de certificados & Mês 2 & Toda a Equipe \\
5. Avaliação & Tabulação de dados, relatórios e prestação de contas & Mês 3 & Coordenação \\
\bottomrule
\end{tabularx}
\end{table}

\vspace{0.3cm}

% ---------------------------------------------------------
% 7. Planilha Orçamentária Detalhada
% ---------------------------------------------------------
\section{Planilha Orçamentária Detalhada}

Todos os custos do projeto são orçados com base em valores médios de mercado para execução de atividades do terceiro setor:

\begin{table}[h!]
\centering
\small
\renewcommand{\arraystretch}{1.2}
\begin{tabularx}{\textwidth}{@{}>{\bfseries}c X c r r@{}}
\toprule
Item & Especificação do Recurso / Serviço & Unid. & Quant. & Total (R\$) \\
\midrule
01 & Honorários de Instrutoria Especializada & Serviço & 1 & 1.200,00 \\
02 & Material Didático, Pastas e Impressão & Kit & 30 & 450,00 \\
03 & Lanche Saudável e Hidratação para Participantes & Refeição & 120 & 720,00 \\
04 & Sonorização, Equipamentos e Infraestrutura de Apoio & Diária & 4 & 600,00 \\
05 & Emissão de Certificados e Materiais de Encerramento & Unid. & 30 & 180,00 \\
\midrule
\multicolumn{4}{@{}r}{\textbf{VALOR TOTAL DO PROJETO:}} & \textbf{R\$ 3.150,00} \\
\bottomrule
\end{tabularx}
\end{table}

\vspace{0.3cm}

% ---------------------------------------------------------
% 8. Indicadores de Monitoramento e Avaliação
% ---------------------------------------------------------
\section{Indicadores de Monitoramento e Avaliação}

Para assegurar a eficácia social da proposta, adota-se matriz objetiva de indicadores de processo e de resultado:

\begin{table}[h!]
\centering
\small
\renewcommand{\arraystretch}{1.2}
\begin{tabularx}{\textwidth}{@{}>{\bfseries}l X c@{}}
\toprule
Indicador & Meta Estabelecida & Meio de Comprovação \\
\midrule
Frequência & Presença superior a 80\% nos encontros & Lista de presença diária \\
Satisfação & Avaliação positiva superior a 90\% & Questionário das famílias \\
Conclusão & Certificação dos concluintes regulares & Livro de registro e entrega \\
Impacto Comunitário & Presença de público na culminância final & Registro fotográfico e ata \\
\bottomrule
\end{tabularx}
\end{table}

\vspace{0.3cm}

% ---------------------------------------------------------
% 9. Termo de Encerramento e Assinaturas
% ---------------------------------------------------------
\section{Termo de Encerramento e Assinaturas}

Declara-se, para todos os efeitos jurídicos e institucionais cabíveis, a exatidão das informações, metodologias e valores discriminados neste plano de trabalho, comprometendo-se a proponente com a plena aplicação dos recursos e execução das atividades descritas.

\vspace{1.5cm}

\begin{center}
    Guapó - GO, \today.
\end{center}

\vspace{1.5cm}

\noindent
\begin{minipage}[t]{0.48\textwidth}
    \centering
    \rule{0.9\textwidth}{0.5pt}\\[0.1cm]
    \textbf{Pr. Presidente}\\
    Igreja Batista Nacional da Paz de Guapó\\
    Representante Legal
\end{minipage}
\hfill
\begin{minipage}[t]{0.48\textwidth}
    \centering
    \rule{0.9\textwidth}{0.5pt}\\[0.1cm]
    \textbf{Coordenação Geral do Projeto}\\
    Equipe de Projetos de Impacto Social\\
    IBNP Guapó
\end{minipage}

\end{document}
